\chapter*{Conclusiones}
\phantomsection
\addcontentsline{toc}{chapter}{Conclusiones}

A partir de la investigación desarrollada sobre la automatización de la rotación de circuitos eléctricos ante déficit de generación, se arribó a las siguientes conclusiones:

\begin{enumerate}
    \item \textbf{Metodología ágil y fundamentos de automatización.} La adopción de la Programación Extrema (XP) permitió gestionar el desarrollo de forma iterativa, adaptándose a los requisitos cambiantes del Despacho Provincial de Carga. La integración de fundamentos de sistemas de gestión energética (EMS) y la virtualización mediante arquitecturas web reactivas justificó la automatización del proceso manual, transformando la infraestructura del Despacho en una plataforma digital escalable y moderna basada en microservicios.
    
    \item \textbf{Diagnóstico y automatización del proceso.} El diagnóstico inicial evidenció que el proceso manual consumía entre 30 y 45 minutos por evento, con errores de cálculo del 5–10\,\% y falta de estandarización. La solución implementada con Next.js y React automatiza el cálculo del déficit mediante un algoritmo FIFO basado en tiempo de estado y proporciona una interfaz gráfica reactiva con componentes especializados (\texttt{RotacionModal}, \texttt{Dashboard}, \texttt{CircuitosTable}), eliminando consultas manuales y reorientando el rol del especialista hacia la supervisión de propuestas generadas automáticamente.
    
    \item \textbf{Validación y resultados.} La verificación técnica confirmó la robustez del API Gateway en NestJS (40–60 ms de respuesta) y la correcta gestión del estado con Context API. La validación con escenarios históricos de 2024 demostró una coincidencia del 100\,\% entre las propuestas generadas y las decisiones manuales, reduciendo el tiempo de respuesta a 2–3 segundos (mejora >99\,\%) y eliminando los errores de cálculo. Las pruebas de usabilidad arrojaron una satisfacción de 4.8/5, destacando el acceso vía navegador, la rapidez de Tailwind CSS y la claridad visual del dashboard.
    
    \item \textbf{Tecnologías y arquitectura de microservicios.} La selección de Next.js/React para el frontend, NestJS con microservicios especializados (circuitos-ms, rotaciones-ms, auth-ms) para el backend y SQL Server 2008 como persistencia legacy respondió a los criterios de accesibilidad multiusuario, compatibilidad con infraestructura existente y alta disponibilidad exigida por el DNC. La arquitectura desacoplada con API Gateway y comunicación NATS garantizan escalabilidad horizontal, mantenibilidad y despliegue centralizado, superando las limitaciones de los sistemas de escritorio. Los microservicios operan sobre tablas maestras de SIGERE (\texttt{ap\_circuitos}, \texttt{ap\_curvas}, \texttt{ap\_apagones}) y PSFV (\texttt{ap\_Aseguramientos}), asegurando integridad referencial y performance optimizado.
    
    \item \textbf{Requisitos y cobertura funcional.} Se formalizaron funcionalidades distribuidas en cinco módulos lógicos: (1) seguridad y autenticación (JWT, roles, Redis blacklist), (2) gestión de circuitos (CRUD, consumo horario, filtrado por bloque/apagable), (3) historial de apagones (consultas optimizadas con CTEs), (4) aseguramientos programados (rango de fechas) y (5) rotación automática (algoritmo FIFO por tiempo de estado, exportación Excel). Los requisitos no funcionales (rendimiento, seguridad, disponibilidad, usabilidad, interoperabilidad, mantenibilidad y escalabilidad) aseguran calidad y sostenibilidad técnica, con tiempos de respuesta <100 ms en backend y <3 s en operaciones completas.
\end{enumerate}